\section{My role}

\textbf{1a.}~At the wk12 kickoff we split responsibilities by module: Harish took hardware and communications, Zian took path planning, Robin took state machines, and Edward and I shared computer vision---I delegated the model training to Edward so I could focus on the detection-to-GPS pipeline. That was the explicit allocation. Table~1 shows where reality diverged.

\vspace{0.3em}
{\small
\noindent\textbf{Table 1: Assigned vs.\ actual responsibilities}

\begin{tabularx}{\textwidth}{@{}l L L@{}}
\toprule
\textbf{Member} & \textbf{Assigned (wk12)} & \textbf{Actual (by wk22)} \\
\midrule
Apollo & CV pipeline (shared w/ Edward) & CV pipeline, simulation framework, ground station, safety layer, progressive test ladder, documentation, complaint letter, report coordination \\
Robin & State machine / flight execution & State machine, integration with CV output, flight-day piloting coordination \\
Harish & Hardware setup, comms between components & Hardware wiring, Cube--Pi serial link, \textcolor{red}{[APOLLO: confirm what else Harish delivered]} \\
Edward & CV model training (delegated by Apollo) & Trained a custom model; we didn't use it---went with standard YOLO person detection instead. \textcolor{red}{[APOLLO: did Edward do anything else?]} \\
Zian & Path planning & Path planning module. \textcolor{red}{[APOLLO: confirm scope---did Zian's planner make it into the final system?]} \\
\bottomrule
\end{tabularx}
}
\vspace{0.3em}

The pattern is visible in the right-hand column: my row kept growing while others stayed roughly static. Every pickup after wk12 was triggered by a gap nobody else was filling---camera calibration because nobody else had the camera, simulation because three cancelled flight days left us with no hardware access, the ground station because the Pi needed headless browser-based control. I initially framed this drift as diligence. I now read it differently. Drawing on Lencioni's (2002) five-dysfunctions model from our Professional Practice sessions, solo building was partly an \emph{absence-of-trust} pattern: I didn't trust the team to navigate work I cared about, so I ring-fenced it. My Ukrainian engineering instinct says one person delivers; this project showed me that instinct is a failure mode dressed as a strength.

The role expansion also had a structural cause I should name honestly: Edward and Zian weren't self-directed. They did what they were asked to do---and that's fair---but they weren't driving things forward on their own initiative. That meant Robin and I absorbed the unassigned work by default, and I absorbed more of it than Robin because I had the laptop-based simulation path while Robin's work depended on hardware access we didn't have.

\textbf{1b.}~The contribution I'm proudest of isn't the 820-plus commits or the retrained model. It's what happened on 2026-03-11, the first cancelled flight day. Weather killed our only window. Instead of going home I calibrated the focal length in the lab, ran lens calibration, wrote a quick-reference sheet for field use without internet, and closed ten commits before dark. That day crystallised my thesis: \emph{the simulation framework was the only reason this project produced verifiable engineering output at all.} Given that we had one real flight video from Sid---one---I built the entire GPS estimation pipeline, calibration tools, and video analysis scripts around it. I squeezed everything I could out of extremely limited data, and I'm genuinely proud of that.

The BGR camera debugging (2026-02-17) is the incident I'd call reflection-in-action (Sch\"on, 1983). I spent hours trying AWB modes, manual gains, ISP tuning---all wrong. Then I stopped guessing and systematically tested all six channel permutations. Root cause: the IMX296 outputs BGR despite labelling itself RGB888. The ratio---three hours of intuition before switching to systematic diagnosis---tells me something about how I respond to frustration. I reach for plausible explanations before exhaustive ones, and the switch happens only after frustration builds up enough to override the comfort of guessing.

\textbf{1c.}~The distinction I'll focus on in future work is between \emph{capability} and \emph{enablement}. The evidence: the \texttt{--robin} flag and HTTP polling interface I built in wk20 let Robin integrate my vision output without touching my code. He worked for two days without asking a single question. That was the first teamwork move I made that scaled without my presence---and it happened in wk20, not wk12. In my next team project I'll schedule explicit pair-programming sessions from week one. Falsifiable signal: can a teammate complete a handoff task independently within 48 hours? If not, my documentation has failed, not their attention.
